\pagenumbering{roman}
\setcounter{page}{1}

\selecthungarian

%----------------------------------------------------------------------------
% Abstract in Hungarian
%----------------------------------------------------------------------------
\chapter*{Kivonat}\addcontentsline{toc}{chapter}{Kivonat}

A szakterület-specifikus nyelvek (DSL) az informatika számos részén jól
használható és kényelmes lehetőséget nyújtanak egy-egy speciális feladatkör
megoldására. Manapság számos olyan eszköz elérhető, amely DSL-ek fejlesztésére
alkalmas, de kevés olyan programnyelv létezik, amelybe a DSL-ek kényelmesen
beágyazhatók.

A dolgozatban egy általam kifejlesztett általános célú programnyelvet mutatok
be, amelyben – a nyelv kialakításának köszönhetően – könnyedén lehet DSL-eket
definiálni és futtatni. A nyelv fordítója LLVM technológiára épít, így akár JIT
segítségével, akár teljes natív fordítással lehetővé teszi a megvalósított
DSL-ben írt programok futtatását.

A dolgozatban demó jelleggel bemutatok néhány gyakorlatban is használható
DSL-t, valamint egy nagyobb léptékű jövőbeli bővítési tervet.

\vfill
\selectenglish


%----------------------------------------------------------------------------
% Abstract in English
%----------------------------------------------------------------------------
\chapter*{Abstract}\addcontentsline{toc}{chapter}{Abstract}

Domain-specific languages (DSLs) are useful in numerous fields of computer
science and provide a convenient way to solve a specific set of tasks. Today
there are many tools available for developing DSLs, but few programming
languages exist into which DSLs can be easily embedded.

In this paper, I present a general-purpose programming language that I
developed, in which DSLs can be easily defined and executed thanks to the way
the language is designed. The language's compiler is based on LLVM technology,
so it allows programs written in the implemented DSL to be executed either with
JIT or with full native compilation.

As a demonstration, I present a few DSLs that can be used in practice, as well
as a larger-scale plan for future expansion.

\vfill
\selectthesislanguage

\newcounter{romanPage}
\setcounter{romanPage}{\value{page}}
\stepcounter{romanPage}
